\documentclass{article}

\begin{document}

pag39

\hfill

\textbf{EXEMPLARY
BIFURCATION
SEQUENCES}

\hfill

A bifurcation sequence is a sequential occurence
of intersting bifurcation events in a family of maps,
seen as the parameter is varied.
In the next four chapters we describe four particularly
exemplary sequences. Each is taken from the
recent literature and illustrates a basic type of bifurcation
behavior. Each chapter is devoted to one of
these types, and may begin with an introduction to
the basic concepts in I D and 2D. The four types of
bifurcation are:

\begin{itemize}
\item absorbing areas,
\item holes,
\item fractal boundaries, and
\item contact bifurcations.
\end{itemize}

\textbf{C4 
ABSORBING AREAS}

\hfill

We begin with a brief introduction to the concept of absorption
in one and two dimensions, and then study an exemplary bifurcation
sequence.


\textbf{4.1 ABSORPTION CONCEPTS, I D}

\hfill

We introduced in Chapter 2 the notions of critical points, which
bound zones of multiplicity, and trapping intervals, which are
mapped into themselves. These notions come together in the concept
of an absorbing interval. This is an interval in the domain of
the map which is trapping, is bounded by critical points, and is
super-attracting, which means that every point sufficiently close to
the critical endpoints will jump into the absorbing interval after a
finite number of applications of the map.
In the context of an iterated map of an interval, the interesting
dynamics take place within absorbing intervals. The critical points
of a one-dimensional map determine absorbing intervals, and are
useful in characterizing some bifurcations, especially those called
global bifurcations. Examples of global bifurcations occur in Chapter
7.

\hfill

\textbf{4.2 ABSORPTION CONCEPTS, 2D}

\hfill

In two-dimensional iterations, we have a notion of absorbing
area, generalizing the absorbing intervals of the one-dimensional
case. The critical curves of a map of the plane playa role analogous
to that of the critical points of a one-dimensional map: they are useful
in determining absorbing areas. We will now illustrate the role
of critical curves in determining these important areas in which the
interesting dynamics occur.
In this chapter we will study the first family of quadratic maps
defined in the Introduction. We will now use the map of this family
determined by b = - 0.8 to illustrate the use of critical curves to
determine an absorbing area.

Our map has two fixed points, P and Q. The basic critical
curve, L, is the locus of points with "coincident preimages," that is,
the set of points having nearby points with different numbers of
preimages or rank 1 (see Appendix 3, especially A3.2 - A3.3, for
definitions). The fundamental critical curve $L_{-1}$ is defined as the
preimage of L. Thus, L is the image of $L_{-1}$ under the map. Similarly,
the derived critical curve LI is the image of L, and so on. All
of these are called critical curves, as described in Chapter 2.

Note for those who have studied vector calculus: In the context
of a generic smooth map, the fundamental critical curve $L_{-1}$ will be
a subset of the set of critical points in the Jacobian sense, points at
which the Jacobian derivative of the map (a linear transformation)
is degenerate (not a linear isomorphism), while the basic critical
curve L is a subset of the set of critical values in the Jacobian sense.
Inflection points are Jacobian critical points which do not belong to
$L_{-1}$.

For this particular map, $L_{-1}$ is the vertical axis, x = 0, and L is a
horizontal line, y = b = - 0.8. It will be convenient to choose a
bounded interval in $L_{-1}$, $S-{-1}$ with endpoints $a_{-1}$ which is the origin
(0, 0), and ao' which is the image of $a_{-1}$ under the map, the
point (0, - 0.8).

Note: The critical curve denoted by L in the text is denoted by
Lo in the figures.
Figure 4-1 shows all these features and more. Note the points
$a_{-1}$ and $a_0$ and the segments of L_I ' L, LI ,00', L 4 • As the map
moves L_I to L, and the point a_I to ao' L is moved to LI ' and the
point ao to a point al in LI . The transversal! (and here, indeed,
orthogonal) crossing of L_I and L at ao is transformed into a tangent
contact of L and LI at a l . Then this point is mapped to a

tangency of LI and L2 at a2' and so on. Note that the curve L2
crosses L_I at the point Po. so L3 is tangent to L at the point PI'
the image of Po. and so on. Similarly. the curve L3 crosses L_I at
the point bo• so L4 is tangent to L at bl • the image of bo• and so
on.
It is worthwhile to pause here and carefully study Figure 4-1. A
point of transversal crossing of any curve. C. through L_I is
mapped into a point of tangency of the image of that curve. ftC).
with L. This is because of the folding which occurs as L_I is
mapped onto L. Also. a point of tangency of a curve C to the curve
L_I is mapped into a point of tangency of the image curvef(C) and
L.
The action of the map may be visualized as a nonlinear folding
of the plane at the fundamental critical curve. L_I • followed by a
nonlinear rotation moving L_I to L around the point ao. and a nonlinear
translation horizontally along L. so that a_I ends up at ao .
The critical curves shown in Figure 4-1 are a kind of skeleton of the
map. as we shall see.

Our next goal is to use these critical curves to discover an
absorbing area of the map. By definition (see Appendix A3.5), an
absorbing area is a region of the plane such that:
• it is mapped into ( or onto) itself;
• its boundary is made up of segments of critical curves, or of
Ii,mit points of an infinite sequence of critical curves; and
• it has a neighborhood every point of which eventually
moves into the absorbing area.
As an example, note in Figure 4-1 that the arcs b 1 a 1 of L, a 1 a2
of LI ' a2a3 of L2, a3a4 of L3 ' and a4bl of L4 bound a region d',
shown shaded in Figure 4-1. This shaded region happens to be an
absorbing area! First of all, it is invariant. For example, the segment
a4bl of L4 is on the boundary of the region, d', but the image of
this arc, aSb2 , is internal to d'. This is clearly shown in Figure 4-2,
in which the image of a4bl is indicated in Ls .
Also, d' is absorbing: A point external to d' is mapped into d'
in a finite number of iterations, as shown in Figure 4-3, and this is
the case for all points sufficiently near to d' .

Figure 4-2. also shows a shaded area, d~, bounded by the critical
curve segments introduced in Figure 4-1. It surrounds a hole, W,
in which the fixed point Q is located. This is a smaller absorbing
area, and is called an annular absorbing area for obvious reasons.
Generally, absorbing areas may be topologically more complex,
with many holes, and with many separate pieces.
Absorbing areas must contain attractors. One, d, is shown as a
cloud of dots in Figure 4-3. It is the attracting set of d', and of the
smaller absorbing area, d~, as well.
Usually there are smaller and smaller absorbing areas around
any attractor of the map. All these, by definition, are bounded by
critical arcs. And as they get smaller and smaller, but always
enclose the same attractor, it is to be expected that the attractor
itself is bounded by critical arcs, or by limit points of infinite
sequences of critical arcs. Figure 4-4 shows 25 iterates of two intervals
of L_l ' the two pieces of d (") L_l . These iterates bound the
attractor shown in Figure 4-3 quite closely.
Another basic concept of dynamics is the basin of an attractor.
In the method of critical curves, we usually determine an absorbing
area by the method illustrated above, and thus the attractor within it,
and then determine the basin of attraction of the absorbing area. In
Figure 4-5 the basin of attraction, D( d'), is shown as the white
region. Its boundary is made up of the repelling (nodal) fixed point,
P, the saddle 2-cycle {Ql' Q2}' and its insets. 1 The points of the
gray region go off to infinity. That is, their trajectories are
unbounded. We call this set the basin of infinity, D( 00). The white
area D(d') (excluding the fixed point Q and its rank 1 preimage
Q_1 in Zo) is the basin of the attractor d.
4.3 EXEMPLARY BIFURCATION SEQUENCE
In this first exemplary bifurcation sequence, we use the first
family with a = 0.7, and decrease b from - 0.4 to - 1.0, in seven
stages.2 For these values of b, our map always has two fixed points,
P and Q, given by:
1. By inset of P we mean the set of points attracted to P, also called the stable set of P.
2. We follow the paper BB.

P: x = 2 ' Y = (1 - a )x
(l-a-Jt»
Q: x = 2 ' Y = (1 - a )x
.P
D(d')
Q._ 1
where () = (1 - a)2 - 4b . Also, there is a 2-cyc1e, {Q 1 ' Q2 }.
Stage I: b = - 0.4
Here, the point Q, at about (- 0.5, - 0.15), is an attractive fixed
point. As b decreases, a Neimark-Hopf bifurcation occurs. The
fixed point Q becomes a repellor and the curve r appears as an
attractive invariant cycle, r (a closed curve mapped onto itself) that
gradually increases in size as b continues to decrease.

Stage 2: b = - 0.5
In Figure 4-6, we see the point repellor Q within the attractive
cycle r , and surrounding that, our first absorbing area, d' , shown
shaded in Figure 4-7. This area is mapped into itself, is bounded by
arcs of critical curves, and is attractive (see Appendix 3.5 for the
precise definition). In this case, the absorbing area is bounded by
the arcs of the critical curves, L, Ll ' L 2 , L3 ' and L 4. These bounding
arcs are generated by successive iterations of the map, as we
now describe.
Notice in Figure 4-7, which shows the critical arcs in more
detail, that L_l and L are straight lines, crossing orthogonally in
one point. Let ao denote this point, (0, - 0.5). Since it belongs to L,
it has a unique preimage, a_I' in L_l . In this case a_I is the origin
(0,0).
Keeping the interval a_I ao in mind, we now draw the successive
images of the halfline of L_l issuing from a_I and containing
ao' that is, issuing downwards. The fourth image, lying within L3 '
crosses the interval a_I ao' as shown in Figure 4-7. At this event
our constructive procedure ends, we have found an absorbing area,
shown shaded in this figure. Further successive images of the segment
converge on r, as shown in the blowup, Figure 4-8. This
procedure may be called Procedure 1. A related procedure is illustrated
in the next stage.
As b continues to decrease, r expands further and eventually
crosses L_l .
Stage 3: b = - 0.6
In this case r crosses L_l in the two points Po and qo' as
shown in Figure 4-9. The wavy shape of r is a consequence of this
crossing, and may be understood as follows.
Apply the map to the configuration shown in Figure 4-9. The
points Po and qo are mapped into the points PI and ql ,also on r.
The straight line segment of L_l between Po and qo is mapped
into the straight line segment of L between PI and q 1 • Because the
map folds the plane two-to-one while moving L_l to L, the curved
segment of r between Po and qo is carried into the curved segment
of r between PI and q 1 ' which is above the line L. The

transversal crossings of r through the line L_I are mapped into
tangencies of r with L at the points P I and q I . These tangencies
are shown clearly in the enlargement, Figure 4-10. This is a universal
property of curves crossing L_I : All crossings are mapped into
tangencies, or contacts, because the map folds the plane at L_I and
maps the two sides of L_I onto just one side of L. Hence, r obtains
a wave from the image of its segment which has crossed L_I .
Another effect of these tangencies is that r is now tangent to
the boundary of the absorbing area d' identified in Stage 2 above,
as shown in Figure 4-11. This absorbing area may be found by the
following method, called Procedure 2.
Consider the straight line segment S_I from a_I to ao in L_I
as above, and construct its successive images am am + I by repeated
applications of the map, until the first crossing with L_I' in the
point bo. See that the first image of S_I ' So' is a straight line segment
from ao to a I in L. The second image, S I ' is a wave from a I
to a2 in L I , likewise S2 in L 2, S3 in L 3 , and S4 in L 4 . But S4
crosses L_I ' and thus bo is found. Let A_I denote the straight line
segment from bo to ao in L_I . Then A_I contains the segment S_I
constructed just above, and its image Ao is a straight line segment
from b l to al in L, containing So. Now the curve segments Ao'
S I ' S2' S3' B enclose the absorbing area d', where B is the curve
segment from a 4 b I within S 4 and L 4' as shown in Figure 4-12.
At this stage, we may see yet another absorbing area, which is
annular in shape. That is, it has a hole. This is shown bounded by
shaded curves in the enlargement of Figure 4-13. Its boundary is
constructed of successive images of the straight line segment A_I
from bo to ao in L_I . The attractive invariant curve, r, is tangent
to the external boundary of this annular absorbing area, as well as to
its interior boundary.
As b decreases further, many bifurcations occur in the dynamics
within these absorbing areas. Probably they have not all been discovered,
but we show just a few events in the remaining figures of
this chapter.

Stage 4: b = - 0.72
At this stage there is an attractive periodic cycle of period 11.
The points of this cycle are labelled in iteration sequence in Figure
4-14. This II-cycle persists as b decreases further, through very
many more bifurcations. The movie on the CD-ROM reveals an
astonishing number of these, and many more have been observed,
even in an interval of b values as narrow as 0.001.
Stage 5: b= -0.78
At this stage we find two attractors coexisting within the annular
absorbing area, a 28-cycle and an II-piece chaotic attractor.
These are shown in Figure 4-15. Near b = 0.798, there is an explosion
to a chaotic attractor filling an annular absorbing area.

Stage 6: b = - 0.7989995
Figure 4-16 shows the chaotic attractor, bounded by critical
curves. Passing below b = - 0.8, there are a number of additional
bifurcations which have been studied on the research frontier. Some
of them will be described later in this book. Approaching b = - 1.0,
further explosions are found.
Stage 7: b = - 0.975
The densely dotted region of Figure 4-17 is a large chaotic
attractor, an annular chaotic area, d. The frontier, F, of its basin of
attraction, includes the inset of the 2-cycle, {Ql ' Q2 }. The boundary
of d is very near F. This figure shows that the critical curves (on
the boundary of the former chaotic area) are about to touch (and
then to cross) the inset of the 2-cycle.

The enlargement of Figure 4-18 shows that, in addition, a contact
of the frontier, F, with the boundary on L is about to occur. This
contact bifurcation, described in Chapter 7, has the effect of
destroying the chaotic attractor, or rather, of transforming it into a
chaotic repellor. Now, almost all of the trajectories diverge to infinity,
except for a Cantor set surviving inside the former chaotic area.
A rough idea of the basin of infinity, D( 00 ), is shown in Figure
4-19 as a black area. The basin of infinity includes infinitely many
holes in the former absorbing area, d', only a few of which are
shown in this figure. The light area is the basin of attraction of the
attractor d. An enlargement is shown in Figure 4-20.

\end{document}
